% !TEX program = xelatex
% ---------------------------------------------------------------------------
%  Xiaoxing Hu · Academic CV (English)
%  Build with: xelatex CV-en.tex
% ---------------------------------------------------------------------------
\documentclass[a4paper,10.5pt]{article}

\usepackage[left=1.75cm,right=1.75cm,top=1.6cm,bottom=1.6cm]{geometry}
\usepackage{xcolor}
\usepackage{titlesec}
\usepackage{enumitem}
\usepackage{tabularx}
\usepackage{array}
\usepackage{fontspec}
\usepackage{xeCJK}
\usepackage{needspace}
\usepackage{fontawesome5}
\usepackage[hidelinks]{hyperref}
\hypersetup{linkcolor=black, urlcolor=black, citecolor=black}

% ── Palette ────────────────────────────────────────────────────────────────
\definecolor{accent}{HTML}{000000}
\definecolor{muted}{HTML}{000000}
\definecolor{rule}{HTML}{000000}

% ── Fonts ──────────────────────────────────────────────────────────────────
\setmainfont{Charter}[
  BoldFont       = Charter Bold,
  ItalicFont     = Charter Italic,
  BoldItalicFont = Charter Bold Italic
]
\setsansfont{Helvetica Neue}
\setmonofont{Menlo}
\setCJKmainfont{Songti SC}[BoldFont=Songti SC Bold]
\newCJKfontfamily\xhei{Heiti SC}[BoldFont=Heiti SC, ItalicFont=Heiti SC,
  BoldFeatures={FakeBold=1.8}, ItalicFeatures={FakeSlant=0.2}]

\pagestyle{empty}
\pagecolor{white}
\setlength{\parindent}{0pt}
\linespread{1.15}
\setlength{\parskip}{0pt}

% ── Section headings ───────────────────────────────────────────────────────
\titleformat{\section}
  {\large\bfseries\color{accent}}
  {}{0em}{}
  [\vspace{-0.85em}{\color{rule}\rule{\textwidth}{0.5pt}}]
\titlespacing*{\section}{0pt}{1.1em}{0.55em}
% keep a heading with the content that follows it
\let\oldSection\section
\renewcommand{\section}[1]{\needspace{5\baselineskip}\oldSection{#1}}

% ── Lists ──────────────────────────────────────────────────────────────────
\setlist[itemize]{
  leftmargin=1.35em,
  topsep=2pt,
  itemsep=1.5pt,
  parsep=0pt,
  label={\color{accent}\footnotesize\textbullet}
}

% ── Macros ─────────────────────────────────────────────────────────────────
% \entry{organisation}{role}{dates}{optional note}
\newcommand{\entry}[4]{%
  \begin{tabularx}{\textwidth}{@{}X r@{}}
    {\bfseries #1}\hfill{\itshape\color{muted}#2} & {\small\color{muted}#3}
  \end{tabularx}%
  \par\vspace{1.5pt}%
  \if\relax\detokenize{#4}\relax\else{\small #4}\par\fi
  \vspace{0.55em}%
}

% \pub{title}{authors}{venue}{links}
\newcommand{\pub}[4]{%
  \needspace{4\baselineskip}%
  \vspace{0.32em}%
  {\bfseries #1}\par\vspace{0.5pt}%
  {\small #2}\par\vspace{0.5pt}%
  {\small\itshape\color{accent}#3}\par\vspace{0.5pt}%
  {\footnotesize\color{muted}#4}\par
}

\newcommand{\sep}{\,\textperiodcentered\,}

% ---------------------------------------------------------------------------
\begin{document}

% ── Header ─────────────────────────────────────────────────────────────────
\begin{center}
  {\fontsize{22}{26}\selectfont\bfseries Xiaoxing Hu}\\[6pt]
  {\small
    \faEnvelope\ \href{mailto:xiaoxinghhh@gmail.com}{xiaoxinghhh@gmail.com}\quad
    \faPhone\ 15111480307\quad
    \faWeixin\ 15111480307\quad
    \faGoogle\ \href{https://scholar.google.com/citations?user=zBM8_XkAAAAJ&hl=zh-CN}{Google Scholar}
  }
\end{center}

\vspace{-0.35em}
{\color{rule}\rule{\textwidth}{0.6pt}}

% ── Education ──────────────────────────────────────────────────────────────
\section{Education}

\entry{Shanghai Jiao Tong University}{Ph.D. Student, Control Science and Engineering}{Sep. 2026 -- Present}
  {}

\entry{Shanghai Innovation Institute}{Ph.D. Student (Joint Program)}{Sep. 2026 -- Present}
  {}

% ── Research interests ─────────────────────────────────────────────────────
\section{Research Interests}

{\bfseries Prior Research:} My prior work focused on foundation vision models and their remote sensing applications, cross-modal representation and alignment, and audio-visual understanding and evaluation. Specific contributions include vision--language alignment and multimodal retrieval, evidence-grounded multimodal reasoning benchmarks with automated evaluation pipelines, and transfer learning and domain generalization for vision foundation models.

\vspace{0.45em}
{\bfseries Current Research Interests:} My current focus is on the full-stack of agent systems, spanning agent models, data, system architecture, and large language model post-training (OPD, MOPD, SWE-OPD).

% ── Publications ───────────────────────────────────────────────────────────
\section{Publications}
{\bfseries\footnotesize\color{muted}Bold denotes the author.\quad
 $^{*}$ Equal contribution.\quad $^{\dagger}$ Corresponding author.\quad $^{\ddagger}$ Student first author.\par}
\vspace{0.2em}

\pub{EvaOmni: An Evidence-Grounded Benchmark for Audio-Visual Reasoning}
  {\textbf{Xiaoxing Hu}, Wenxi Yao, Chaoyuan Shi, Ziqi Ye, Xiaorong Guo, Peixuan Li, Zhuoran Sun,
   Ziyang Gong, Peiyuan Zhang, Mingxin Liu, Shuran Ma, Hongbo Liu, Xiaomin Yu, Yu Tian, Wentao Wang,
   Zhenjie Yang, Jian Hu, Yunhang Shen, Haoyu Cao, Xue Yang$^{\dagger}$}
  {Under Review}
  {\href{https://xiaoxing2001.github.io/EvaOmni.github.io/}{Project}\quad
   \href{https://github.com/VisionXLab/EvaOmni}{Code}\quad
   \href{https://arxiv.org/abs/2610.28444}{Paper}}

\pub{ProCLIP: Progressive Vision-Language Alignment via LLM-based Embedder}
  {\textbf{Xiaoxing Hu}, Kaicheng Yang, Ziyang Gong, Qi Ming, Zonghao Guo, Yu Tian, Xiang An,
   Ziyong Feng, Xue Yang$^{\dagger}$}
  {NeurIPS 2026 (Spotlight)\sep CCF-A}
  {\href{https://github.com/VisionXLab/ProCLIP}{Code}\quad\href{https://arxiv.org/abs/2510.18795}{Paper}}

\pub{Earth-Adapter: Bridge the Geospatial Domain Gaps with Mixture of Frequency Adaptation}
  {\textbf{Xiaoxing Hu}$^{*}$, Ziyang Gong$^{*}$, Yupei Wang$^{\dagger}$, Yuru Jia, Fei Lin, Dexiang Gao, Ke An,
   Jianhong Han, Zhuoran Sun, Gen Luo, Xue Yang$^{\dagger}$}
  {AAAI 2026\sep CCF-A}
  {\href{https://github.com/VisionXLab/Earth-Adapter}{Code}\quad
   \href{https://arxiv.org/abs/2504.06220}{Paper}}

\pub{Decoupled Global-Local Alignment for Improving Compositional Understanding}
  {\textbf{Xiaoxing Hu}$^{*}$, Kaicheng Yang$^{*}$, Jun Wang, Haoran Xu, Ziyong Feng, Yupei Wang$^{\dagger}$}
  {ACM MM 2025\sep CCF-A}
  {\href{https://xiaoxing2001.github.io/DeGLA.github.io/}{Project}\quad
   \href{https://github.com/xiaoxing2001/DeGLA}{Code}\quad
   \href{https://arxiv.org/abs/2504.16801}{Paper}}

\pub{Boosting Domain Generalization in Remote Sensing Image Segmentation via Style Mapping and General Prototypical Contrast}
  {Yupei Wang, \textbf{Xiaoxing Hu}$^{\ddagger}$, Yongkang Hu, Zhuoran Sun, Shanghang Zhang, Liang Chen$^{\dagger}$}
  {International Journal of Computer Vision (IJCV) 2025\sep CCF-A}
  {\href{https://link.springer.com/article/10.1007/s11263-025-02568-3}{Paper}}

% ── Internships ────────────────────────────────────────────────────────────
\section{Internships}

\entry{Tencent Youtu Lab \sep ShenNong Research Center}{Research Intern}{Dec. 2025 -- Sep. 2026}
  {\begin{itemize}
     \item Contributed to omni-modal model development, owning the data, training, and evaluation loop.
     \item Contributed to building an office-agent evaluation framework; introduced a Rubric-based LLM-as-a-Judge automated evaluation mechanism for PPT-Agent execution trajectories to support agent capability assessment and iterative verification.
     \item Helped establish dataset evaluation and cleaning workflows; wrote and submitted one paper (under review).
   \end{itemize}}

\entry{DeepGlint \sep Foundation Model Group}{Research Intern}{Jun. 2025 -- Oct. 2025}
  {\begin{itemize}
     \item Conducted research on multimodal retrieval foundation models, focusing on vision--language representation alignment.
     \item The resulting first-authored work was accepted to NeurIPS 2026 as a Spotlight (ProCLIP).
   \end{itemize}}

% ── Honors ─────────────────────────────────────────────────────────────────
\section{Honors \& Awards}

\begin{itemize}
  \item \textbf{Xu Teli Scholarship Candidate} --- 6 graduate students selected university-wide\hfill Jun. 2026
  \item \textbf{National Scholarship}\hfill Oct. 2025
\end{itemize}

% ── Skills ─────────────────────────────────────────────────────────────────
\section{Skills}

\begin{itemize}
  \item \textbf{Languages}: Python (proficient), C/C++, Shell
  \item \textbf{Tools}: PyTorch, Hugging Face Transformers, DeepSpeed, Git, \LaTeX
  \item \textbf{Domain}: multimodal pretraining and alignment, instruction tuning, retrieval, benchmark construction
\end{itemize}

\end{document}
